\section{Security}

We now present proofs of
correctness and a high-level overview of our proof of security against
chosen-message attacks for an adaptive version of FROST,
with a heuristic assumption on how the proof of security applies to FROST directly.

Note that subsequent works have given a direct proof for FROST~\cite{}.

\subsection{Correctness}
\label{frost:correctness}

Signatures in FROST
are constructed from two polynomials;
the first polynomial $F_1(x)$ defines the secret sharing of the private
signing key $s$ (such that $Y = g^s$) and the
second polynomial $F_2(x)$ defines the secret sharing of the nonce $k$
such that
\(
k = \sum_{i \in S} d_i + e_i \cdot \rho_i
\)
using the associated public data $(m, B)$ to determine $\rho_i$.
During the key generation phase described in Figure~\ref{fg:keygen},
the first polynomial $F_1(x) = \sum_{j=1}^n f_j(x)$ is generated
such that the secret key shares are $s_i = F_1(i)$ and the secret key is $s = F_1(0)$.

During the signature phase (Figure~\ref{fg:single-sign-parallel}),
each of the
participants in the coalition $\coalition$ selected for signing use a pair of nonces $(d_i, e_i)$ to
define a degree $\lvert\coalition-1 \rvert$ polynomial $F_2(x)$, interpolating the
values $(i,\frac{d_i + e_i \cdot H_1(i, m, B)}{\lambda_i})$,
such that $F_2(0) = \sum_{i \in S} d_i + e_i \cdot \rho_i$.

Then let $F_3(x) = F_2(x) + c \cdot F_1(x)$, where
$c = H_2(R, Y, m)$.  Now $z_i$
equals $d_i + (e_i \cdot \rho_i) + \lambda_i \cdot s_i \cdot c =
\lambda_i (F_2(i) + c \cdot F_1(i)) = \lambda_i F_3(i)$, so $z = \sum_{i \in S} z_i$
is simply the Lagrange interpolation of $F_3(0) = (\sum_{i \in S} d_i + e_{ij} \cdot \rho_i) + c
\cdot s$.  Because $R = g^{\sum_{i \in S} d_i + e_i \cdot \rho_i}$ and $c
= H_2(R, Y, m)$, $(R, z)$ is a
correct Schnorr signature on $m$.

\subsection{Security Against Chosen Message Attacks}
\label{frost:high-level-proofs}

We now present a high-level argument of security for FROST.
We begin by summarizing a proof of security for an interactive variant of FROST
that we call FROST-Interactive, and then demonstrate how the proof extends to
plain FROST.

We employ the generalized forking strategy used by Bellare and
Neven~\cite{generalized-forking-lemma} to create a reduction to the security of
the discrete logarithm problem (DLP) in $\mathbb{G}$.
We prove security against the standard notion of existential unforgeability
against chosen message
attacks (EUF-CMA) by demonstrating that the difficulty to an adversary
to forge FROST-Interactive signatures by
performing an adaptively chosen message attack in the random oracle model
reduces to the difficulty of computing the discrete logarithm of an arbitrary
challenge value $\omega$ in the underlying group,
so long as the adversary controls fewer than the threshold $t$
participants.

\textbf{FROST-Interactive.}
In FROST-Interactive, $\rho_i$ is established using a
``one-time'' verifiable random function (VRF),\footnote{A one-time VRF $F_k$ for key $k$ relaxes the standard properties of a
VRF by requiring that $F_k(x)$ be unpredictable to someone who does not know
$k$ only when at most one value of $F_k(y)$ has been published by the keyholder
(and $y \not= x$). We use the construction $k=(a,b) \in \mathbb{Z}_q^2$ and
$F_k(x) = a + b\cdot x$. The public key is $(A=g^a,B=g^b)$.}
as $\rho_i = a_{ij} + (b_{ij} \cdot H_{\rho}(m, B))$,
where $(a_{ij}, b_{ij})$ are selected and committed to as
$(A_{ij}=g^{a_{ij}}, B_{ij}=g^{b_{ij}})$ during the preprocessing stage, along
with zero-knowledge proofs of knowledge of $(a_{ij}, b_{ij})$.
To perform a signing operation, participants first generate
$\rho_i$ in the first round of the signing protocol using $(a_{ij}, b_{ij})$, and then
publish $\rho_i$ to the signature aggregator, which distributes all $\rho_\ell, \ell
\in S$ to all signing participants. These $\rho_\ell, \ell \in S$ values are then
used by all signing participants to compute $R$ in the second round of the
signing protocol, which participants use to calculate and publish $z_i$.

\textbf{Summary of proof for EUF-CMA security for FROST-Interactive.}
Let
$n_h$ be the number of queries made to the random oracle,
$n_p$ be the number of allowed preprocess queries, and
$n_s$ be the number of allowed signing queries.
We assume there exists a
forger $\mathcal{F}$ that
$(\tau, n_h, n_p, n_s, \epsilon)$-breaks FROST-Interactive, meaning that
$\mathcal{F}$ can compute a forgery for a signature generated by
FROST-Interactive in time $\tau$ with success $\epsilon$, but is limited to
making
$n_h$ number of random oracle queries,
$n_p$ number of preprocess queries,
and
$n_s$ number of signing queries.
We construct an algorithm $C$ that $(\tau', \epsilon')$-solves the discrete
logarithm problem in $\mathcal{G}$, for an arbitrary challenge value $\omega
\in \mathbb{G}$,
using as a subroutine a forger $\mathcal{F}$ that can forge FROST-Interactive signatures.

